%% =========================================================
%%  LocalDemand -- Literature Review (IEEE Conference Format)
%%  B.Tech Capstone Project
%%  Sections completed: 1, 2 (of 8)
%% =========================================================

\documentclass[conference]{IEEEtran}
\IEEEoverridecommandlockouts

\usepackage{cite}
\usepackage{amsmath,amssymb,amsfonts}
\usepackage{graphicx}
\usepackage{textcomp}
\usepackage{xcolor}
\usepackage[utf8]{inputenc}
\usepackage[T1]{fontenc}
\usepackage{hyperref}

\begin{document}

\title{LocalDemand: A Literature Review on Demand Forecasting,
Inventory Optimization, and AI Adoption for Food-Sector SMEs}

\author{%
  \IEEEauthorblockN{B.Tech Capstone Project Team}
  \IEEEauthorblockA{\textit{LocalDemand} \\ }
}

\maketitle

% ----------------------------------------------------------
\begin{abstract}
This paper presents a structured literature review of 25 peer-reviewed
research works that collectively underpin the design and rationale of
LocalDemand --- a Flutter-based mobile application that employs
Facebook Prophet to generate daily demand forecasts and actionable
restock recommendations for small food businesses (cafés, bakeries, and
restaurants). The review is organized into eight thematic sections
covering: forecasting methods and time-series models; the integration of
external regressors such as weather and holidays; challenges specific to
SME contexts including sparse data and cold-start; inventory optimization
and waste reduction; technology adoption and accessibility; and a final
synthesis of research gaps addressed by LocalDemand. All findings
reported here are directly drawn from the cited works; no results have
been fabricated or estimated.
\end{abstract}

\begin{IEEEkeywords}
Demand Forecasting, Facebook Prophet, SME, Inventory Management,
Time-Series, LSTM, XGBoost, Mobile Application, LocalDemand
\end{IEEEkeywords}

% ==========================================================
\section{Introduction}
% ==========================================================

Accurate demand forecasting is a cornerstone of efficient operations in
the food-retail sector. For small and medium enterprises (SMEs) such as
independent cafes, bakeries, and neighbourhood restaurants, the ability
to predict daily sales volumes has direct consequences on food waste,
procurement costs, customer satisfaction, and overall profitability.
Unlike large retail chains that can absorb the financial impact of
over-purchasing or stockouts, micro-SMEs operate on thin margins where a
single week of poor inventory decisions can translate into measurable
losses \cite{rahmatillah2026, groene_zakharov}.

The challenges facing SME demand forecasting are multi-dimensional.
Demand in food businesses exhibits strong daily and weekly seasonality,
is sensitive to local weather conditions, fluctuates sharply around
public holidays and local events, and is further complicated by
promotional activities \cite{guler, lalitha2025, jain2024, qureshi2024}.
At the operational level, these businesses generate sparse, sometimes
zero-inflated sales histories that violate the assumptions of classical
statistical models \cite{susanto2026, munawar2026}. On the adoption
side, enterprise-grade AI and ERP solutions are economically and
technically out of reach for most micro-SMEs, which lack both the IT
budget and the data-science expertise required to deploy them
\cite{aree2025, yuniarty2025, srivastava2026}.

Despite a growing body of research on forecasting algorithms---ranging
from classical ARIMA variants to deep learning architectures such as
LSTM and GRU---a persistent gap remains between \emph{forecasting
accuracy} and \emph{operational actionability}. Most studies terminate
at the computation of error metrics (RMSE, MAPE, MAE) without
translating predictions into concrete inventory decisions such as
purchase quantities or restock schedules \cite{aqila2026, munawar2026,
rahmatillah2026}. Furthermore, almost no reviewed work combines
automated forecasting, external regressor integration (weather,
holidays), and a lightweight mobile delivery layer accessible to
non-technical SME owners in a single system.

This literature review is structured to trace the evidence base for each
design decision in LocalDemand. Section~II surveys forecasting
algorithms. Section~III examines the role of external factors.
Section~IV addresses SME-specific data challenges. Section~V covers
inventory optimization and waste reduction. Section~VI analyses
technology adoption barriers and mobile solutions. Section~VII
synthesizes the overall research landscape and explicitly maps each
identified gap to the corresponding design choice in LocalDemand.
Section~VIII concludes.

% ==========================================================
\section{Forecasting Methods and Time-Series Models}
% ==========================================================

The literature on demand forecasting encompasses a broad spectrum of
methodologies, from classical statistical models to modern deep learning
architectures. Understanding the trade-offs among these approaches is
essential for selecting a method that is both accurate and deployable
in a resource-constrained SME environment.

% ----------------------------------------------------------
\subsection{Traditional and Statistical Models}
% ----------------------------------------------------------

Classical time-series methods --- ARIMA, SARIMA, Holt-Winters, and
moving averages --- form the historical backbone of demand forecasting.
Gartner et al.\ (2024) conducted a systematic evaluation of automated
demand forecasting for SMEs, benchmarking classical models against deep
learning methods. Their results showed that traditional statistical models 
were consistently outperformed by machine learning; the best-performing 
Dilated CNN achieved a normalised RMSE of 0.3119, while ARIMA variants
lagged significantly. Nevertheless, the study concluded that classical 
models retain value when historical data is limited, requiring far fewer 
training samples than deep learning \cite{gartner2024}.

Rahmatillah et al.\ (2026) provided a direct contrast between a 7-day 
Moving Average (MA7) and an optimised LSTM for an ice cream SME. The 
moving average, while computationally trivial, failed to anticipate demand 
spikes and led to persistent stockouts in purchasing simulations. This 
finding highlights why LocalDemand avoids moving-average baselines for its 
restock logic, despite their popularity among small business owners 
\cite{rahmatillah2026}.

% ----------------------------------------------------------
\subsection{Machine Learning and Ensemble Models}
% ----------------------------------------------------------

Ensemble and gradient-boosting methods have emerged as robust alternatives
to both classical statistics and computationally expensive deep learning. 
Kamboj et al.\ (2025) demonstrated the power of Facebook Prophet combined 
with a gradient-boosting residual corrector. Their framework applied Prophet 
to decompose trend and seasonality in drugstore sales data, then used XGBoost 
on the residuals to capture nonlinear promotional effects. This hybrid model 
reduced the median MAPE by 12\% to 18\% compared to classical baselines, 
providing compelling empirical validation for Prophet-based architectures 
\cite{kamboj2025}.

Jatte et al.\ (2025) presented a comparative study across six models on a 
retail demand dataset of over 100{,}000 records. While LSTM achieved the 
highest accuracy of 92.31\%, Prophet reached a competitive 85.71\%. This 
benchmark confirms that Prophet delivers strong performance while requiring 
significantly less feature engineering overhead than deep learning 
\cite{jatte2025}. 

Additionally, Groene and Zakharov deployed an XGBoost model on hourly POS 
sales data from a fast-food restaurant. By incorporating weather conditions, 
holidays, and events, their model significantly outperformed manual managerial 
forecasting, reducing RMSE by 22\%--33\%. This demonstrates that ensemble ML 
models with engineered external features can yield quantifiable accuracy gains 
in a single-location SME setting \cite{groene_zakharov}.

% ----------------------------------------------------------
\subsection{Deep Learning Models}
% ----------------------------------------------------------

Deep learning architectures demonstrate peak accuracy on sequential demand 
data, though typically demanding extensive computational resources. Aqila 
et al.\ (2026) evaluated SARIMA, FB Prophet, and a Stacked LSTM model on 
hourly coffee shop POS data. The LSTM achieved the lowest RMSE of 0.1535, 
compared to 0.46 for Prophet. However, the study stopped at error metric 
computation without converting predictions into inventory decisions --- the 
forecast-to-action gap that LocalDemand addresses \cite{aqila2026}.

Munawar and Sudirman (2026) applied an LSTM model to hourly culinary SME 
sales, capturing complex intra-day patterns to achieve the lowest MAE and RMSE. 
While their cost-based simulations showed that LSTM predictions reduced 
overstocking penalties, their model explicitly excluded external drivers like 
weather and holidays, limiting its real-world applicability for seasonally 
sensitive food businesses \cite{munawar2026}.

At the extreme end of complexity, the AID-PAF framework integrated a 
Temporal-Fusion Neural Architecture (TFNA) with waste-aware linear programming. 
Evaluated on a digital twin simulator, it achieved a 6.5\% MAPE and a 40\% 
reduction in food waste. While demonstrating the theoretical ceiling of 
deep learning optimization, its dependence on large labelled datasets and 
high-performance infrastructure makes it unsuitable for resource-constrained 
micro-SMEs \cite{aidpaf}.

% ----------------------------------------------------------
\subsection{Hybrid and Prophet-Based Approaches}
% ----------------------------------------------------------

Facebook Prophet is a decomposable time-series model that estimates trend, 
seasonality, and holiday effects while natively supporting external regressors. 
Its robustness to missing data, ease of configuration, and interpretability 
make it highly attractive for SMEs.

Guler et al.\ applied Prophet directly to restaurant sales forecasting using 
eight months of POS data augmented with weather and holiday markers. The 
model generated 30-day forecasts (MAE = 0.76, RMSE = 0.82), confirming 
Prophet's ability to capture seasonal trends and holiday-driven demand spikes 
in a single-location setting \cite{guler}.

As noted previously, Kamboj et al.\ (2025) successfully used Prophet as the 
decomposition backbone of their hybrid XGBoost model, reducing MAPE by up to 
18\% over classical baselines \cite{kamboj2025}. Similarly, Jatte et al.\ 
(2025) confirmed Prophet's viability with an 85.71\% accuracy on retail data 
--- slightly trailing LSTM (92.31\%) but requiring vastly less data engineering 
and computational overhead \cite{jatte2025}.

Together, this evidence establishes a clear hierarchy: deep learning delivers 
peak accuracy for data-rich environments, while Prophet-based approaches 
provide a pragmatic, highly resilient middle ground for data-sparse SME 
contexts. This validates LocalDemand's choice of Facebook Prophet.

% ==========================================================
\section{Integration of External Factors}
% ==========================================================

A recurring finding across the reviewed literature is that demand for
food and beverage products is not driven solely by historical purchase
patterns. External environmental and contextual factors --- most
prominently weather conditions, public holidays, local events, and
promotions --- account for a substantial share of demand variability.
Ignoring these factors leads to systematic forecast bias, particularly
in small food businesses where the customer base is local and therefore
highly sensitive to immediate environmental conditions.

% ----------------------------------------------------------
\subsection{Weather Conditions}
% ----------------------------------------------------------

Weather is a consistently cited external driver of food demand. Lalitha 
et al.\ (2025) integrated weather and cost data into a Random Forest 
framework to forecast restaurant demand, achieving an MSE of 53.97. They 
used this to develop a dynamic pricing engine, moving beyond pure prediction 
into an actionable revenue response, though their reliance on continuous data 
feeds limits micro-SME applicability \cite{lalitha2025}. 

Jain et al.\ (2024) found that up to 40\% of predicted sales in the food 
sector fail to materialise due to unforeseen external factors like weather. 
Their weather-integrated ML models successfully generated adaptive menu 
recommendations, although stopping short of quantitative restocking outputs 
\cite{jain2024}. 

Furthermore, Groene and Zakharov demonstrated that adding weather, holiday, 
and event data to an XGBoost model reduced forecasting RMSE by 22\%--33\% for 
a fast-food restaurant \cite{groene_zakharov}. These studies collectively 
confirm that weather data provides independent, highly valuable predictive 
signals for food retail.

% ----------------------------------------------------------
\subsection{Holidays, Special Days, and Events}
% ----------------------------------------------------------

Public holidays and special calendar days are universally recognised as 
demand disruptors. Guler et al.\ explicitly incorporated holiday data into 
their Prophet model for restaurant sales. Their results demonstrated that 
Prophet's native holiday component effectively captured demand spikes, 
establishing the Prophet + weather + holidays configuration as practically 
viable for single-location food businesses \cite{guler}.

Kamboj et al.\ (2025) leveraged Prophet's holiday handling specifically so 
their residual learner would not be distorted by calendar anomalies 
\cite{kamboj2025}. In contrast, Gartner et al.\ (2024) noted that their SME 
forecasting system lacked holiday integration, which contributed to higher 
error rates relative to optimized models \cite{gartner2024}. This strongly 
motivates the explicit holiday-regressor design in LocalDemand.

% ----------------------------------------------------------
\subsection{Promotions and Other Regressors}
% ----------------------------------------------------------

Beyond weather and holidays, promotional activities and macroeconomic 
indicators carry predictive value. Qureshi et al.\ (2024) observed that 
including promotions alongside weather yielded better performance than any 
individual subset \cite{qureshi2024}. The AID-PAF framework even successfully 
incorporated social media sentiment as an additional demand signal \cite{aidpaf}. 

In synthesis, weather and holiday regressors are the most consistently 
evidence-backed external demand drivers for food SMEs. Because they are 
publicly available and natively supported by Prophet's regressor framework, 
they form the primary external feature set for LocalDemand, consistent with 
the methodologies of Guler et al., Kamboj et al., and Groene \& Zakharov.

% ==========================================================
\section{Handling Challenges in SME Contexts}
% ==========================================================

Forecasting for micro-SMEs presents a set of challenges that are
qualitatively different from those encountered in large retail
deployments. The reviewed literature identifies three principal
structural challenges: sparse and zero-inflated transaction data,
limited historical depth (cold-start), and the absence of technical
infrastructure to support complex model maintenance.

% ----------------------------------------------------------
\subsection{Sparse and Zero-Inflated Data}
% ----------------------------------------------------------

Many micro-SMEs operate on a pre-order or low-frequency sales model,
generating transaction records that contain a high proportion of
zero-value days. Susanto et al.\ (2026) investigated this phenomenon
directly by analysing sales prediction for an SME operating under a
pre-order framework. The authors benchmarked an XGBoost regression model
against a simple 7-day Moving Average (MA(7)) on this zero-inflated
dataset. Despite its algorithmic sophistication, XGBoost struggled to
accurately model the discontinuous, spike-dominated demand pattern,
with the baseline moving average providing more stable estimates on
several evaluation windows. The study concluded that continuous
regression models are fundamentally ill-suited to zero-inflated,
low-volume data environments, explicitly pointing to the need for models
that can handle intermittent demand more gracefully \cite{susanto2026}.

This finding is directly consequential for LocalDemand. Facebook Prophet
treats missing values natively and does not require a complete, unbroken
time series --- it fits a piecewise-linear or logistic trend to whatever
data points are present and uses Fourier series to estimate seasonality,
making it substantially more robust to zero-inflated and sparse inputs
than regression-tree methods such as XGBoost \cite{susanto2026}.

% ----------------------------------------------------------
\subsection{Limited Historical Data and Cold-Start Problems}
% ----------------------------------------------------------

Even for SMEs with regular daily sales, the total historical depth
available at system launch is typically short --- often less than one
year. Rahmatillah et al.\ (2026) demonstrated that their optimised LSTM
model achieved strong performance (RMSE = 1.65, MAPE = 45.06\%) on an
ice cream SME dataset containing nearly 1{,}000 days of sales history for
a single product. The study's inventory simulation results were positive,
eliminating lost demand entirely under a weekly purchasing cycle.
However, the authors acknowledged that this performance was contingent
on the richness of the training dataset; the model's cold-start
behaviour on shorter histories was not evaluated. This represents a
significant limitation for newly established food businesses or new menu
items, for which no historical precedent exists \cite{rahmatillah2026}.

Gartner et al.\ (2024) encountered similar cold-start constraints in
their automated SME forecasting system. Their evaluation found that the
Dilated CNN and other deep learning models required substantially more
historical data to converge than ARIMA or GAM, making them unreliable
for SMEs with fewer than six to twelve months of records. The study
recommended statistical or Prophet-based models as pragmatic fallbacks
for data-scarce businesses \cite{gartner2024}.

Munawar and Sudirman (2026) also highlighted cold-start limitations in
their LSTM-based hourly forecasting work, noting that the model's
performance degraded significantly when the training window was
shortened below approximately 180 days of hourly records. This finding
reinforces the case for Facebook Prophet, which has been demonstrated in
multiple studies to produce reasonable forecasts with as few as a few
weeks of daily observations, owing to its Bayesian prior structure and
regularised trend estimation \cite{munawar2026}.

% ----------------------------------------------------------
\subsection{Data Scarcity Solutions and Design Implications}
% ----------------------------------------------------------

Across the reviewed works, three categories of solution emerge for
managing SME data scarcity. First, model selection: choosing algorithms
that are inherently robust to short histories and missing values, as
argued by Susanto et al.\ and Gartner et al.\ in favour of statistical
and Prophet-based methods over deep learning \cite{susanto2026,
gartner2024}. Second, external regressor augmentation: supplementing
sparse internal sales data with external signals (weather, holidays)
that carry independent predictive power, as demonstrated by Groene \&
Zakharov, Guler et al., and Kamboj et al.\ \cite{groene_zakharov, guler,
kamboj2025}. Third, simulation and safety-stock buffers: using
conservative purchasing simulations with built-in safety margins to
absorb the uncertainty from short forecast horizons, as practised by
Rahmatillah et al.\ and Munawar \& Sudirman \cite{rahmatillah2026,
munawar2026}.

LocalDemand applies all three strategies simultaneously. It selects
Facebook Prophet as its core model for cold-start resilience; it
integrates weather and holiday regressors as external demand signals;
and it converts Prophet's point forecasts into restock recommendations
that incorporate a configurable safety-stock buffer, ensuring that the
limited historical depth available at system launch does not translate
directly into under-ordering decisions.

% ==========================================================
\section{Inventory Optimization, Waste Reduction, and Actionable Insights}
% ==========================================================

A central shortcoming identified throughout the forecasting literature is
the failure to bridge the gap between prediction and operational action.
Computing an accurate RMSE is not, by itself, useful to a small business
owner deciding how many loaves of bread to order tomorrow. This section
reviews works that take the additional step of converting demand forecasts
into concrete inventory decisions, waste-reduction mechanisms, or
restocking simulations.

% ----------------------------------------------------------
\subsection{From Forecasting to Restock Recommendations and Simulations}
% ----------------------------------------------------------

Rahmatillah, Herwanto, and Sudirman (2026) are among the few authors to
directly simulate the operational consequences of forecast accuracy. After
comparing an optimised LSTM against a 7-day Moving Average (MA7) on daily
ice cream SME sales data (nearly 1{,}000 days of records), the authors
applied each model's outputs to a simulated weekly purchasing cycle. The
optimised LSTM achieved RMSE = 1.65 and MAPE = 45.06\%, outperforming MA7
on error metrics, but the more important finding was in simulation: the
LSTM-driven purchasing policy eliminated lost demand (stockout units) entirely
while generating a negligible overstock of only 1.57 units. The moving
average, by contrast, produced recurring stockouts. This study establishes
the methodological precedent for LocalDemand's restock recommendation
engine: translating a Prophet-generated point forecast into a weekly
purchase quantity with a safety buffer, mirroring the simulation logic
demonstrated here \cite{rahmatillah2026}.

Munawar and Sudirman (2026) similarly extended their LSTM hourly sales
forecasting analysis to include cost-based simulations. By assigning
per-unit costs to overstocking (spoilage) and understocking (lost revenue),
the authors demonstrated that the LSTM model's lower MAE and RMSE
translated directly into reduced total inventory costs compared to
baseline models. This work reinforces that forecast accuracy improvements
are not merely academic --- they carry measurable financial value for SME
operators \cite{munawar2026}.

The Demand Forecasting System for Small Business Optimization study
provided a web interface for uploading sales data and generating DeepAR
forecasts, with outputs intended to support stock planning decisions.
DeepAR achieved 81.0\% forecasting accuracy versus SARIMA at 50.76\% and
a hybrid approach at 77.14\%. While this represents a positive step toward
accessible forecasting tools, the system did not incorporate external
factors (weather, holidays) and its web-based delivery mechanism requires
internet access and a desktop interface --- barriers for mobile-first
micro-SME operators \cite{smallbiz_deepar}.

% ----------------------------------------------------------
\subsection{Sustainability and Food Waste Minimization}
% ----------------------------------------------------------

Food waste reduction is both an economic and an environmental imperative
in the food-service sector. The AID-PAF framework addressed this directly
by coupling a Temporal-Fusion Neural Architecture (TFNA) for demand
forecasting with a waste-aware linear programming (LP) model for inventory
optimisation. The LP layer explicitly penalises over-ordering of perishable
items, balancing the trade-off between customer fill rate and spoilage cost.
Evaluated on an AnyLogic-MATLAB digital twin simulator, the integrated
system achieved a MAPE of 6.5\%, a 96.2\% customer fill rate, a 40\%
reduction in food waste, and a 21\% reduction in procurement costs
\cite{aidpaf}. These results represent the highest reported food-waste
reduction in the reviewed literature and validate the principle that
coupling demand forecasting with a waste-aware inventory optimisation
layer produces substantially better outcomes than forecasting alone.

The Predictive Analytics for Food Supply Chain and Waste study evaluated
LSTM, XGBoost, and Random Forest on restaurant POS data, inventory logs,
and waste records collected over three months. XGBoost achieved the best
performance across all metrics: 94.3\% accuracy, 95.1\% fill rate, 53.8\%
food waste reduction, and 27.6\% cost savings --- outperforming LSTM and
Random Forest. The study demonstrates that even without a dedicated
optimisation layer, accurate demand forecasting significantly reduces waste
and improves profitability \cite{foodwaste_ml}.

Lalitha et al.\ (2025) framed the problem differently: rather than
minimising post-purchase waste, their dynamic pricing engine sought to
reduce the probability of unsold perishable inventory by adjusting prices
in real-time as weather conditions changed. This approach --- preventing
over-purchasing in the first place by modulating demand through pricing ---
complements the inventory-side approach of LocalDemand \cite{lalitha2025}.

% ----------------------------------------------------------
\subsection{Smart Inventory Systems for SMEs}
% ----------------------------------------------------------

The Smart Stock study proposed a hardware-software system for small
warehouse inventory management, integrating XGBoost-based demand
forecasting with QR code product identification, Raspberry Pi edge
devices, IoT sensors, and cloud storage. The proposed system achieved
96\% classification accuracy and an R$^{2}$ value greater than 0.95,
outperforming Random Forest, Linear Regression, and KNN on the same
dataset. Real-time stock monitoring and automated replenishment alerts
helped reduce stockouts, overstocking, and manual data-entry errors,
positioning the system as a cost-effective alternative to enterprise ERP
platforms such as SAP and Oracle \cite{smartstock}.

While the hardware component (Raspberry Pi edge devices, QR readers)
is beyond the scope of LocalDemand, the study's core insight is highly
relevant: SMEs benefit measurably from automated, low-cost inventory
monitoring and decision support. LocalDemand replicates this automation
and accessibility at the pure software level, eliminating hardware
dependencies by delivering forecast-driven restock recommendations through
a mobile application that any SME owner can use with a smartphone.

The synthesis across Section~V reveals a consistent pattern: the most
impactful systems in the reviewed literature are those that do not stop
at generating a forecast but instead translate predictions into a concrete
operational output --- a purchase order, a restock quantity, a pricing
adjustment, or an automated alert. LocalDemand is explicitly designed
around this principle, with its core deliverable being a daily restock
recommendation rather than a raw time-series forecast.

% ==========================================================
\section{Technology Adoption and Accessibility for SMEs}
% ==========================================================

Even the most accurate forecasting model provides no value if SME
operators are unwilling or unable to adopt it. A significant body of
literature examines the behavioural, economic, and technical barriers
that prevent small business owners from integrating AI and ML tools into
their operations, and a smaller but growing set of works proposes
lightweight, mobile-first solutions to address these barriers.

% ----------------------------------------------------------
\subsection{Barriers to AI Adoption: TAM and PLS-SEM Evidence}
% ----------------------------------------------------------

Yuniarty et al.\ (2025) applied the Technology Acceptance Model (TAM)
to survey 131 SME owners in Indonesia about their intention to adopt AI
tools. Using Partial Least Squares -- Structural Equation Modelling
(PLS-SEM), the study found that \emph{perceived ease of use} and
\emph{perceived security} were the primary determinants of adoption
intention. Critically, \emph{perceived usefulness} --- the actual
functional capability of the AI tool --- did not significantly drive
adoption if the system was considered difficult to use or insecure. This
finding has direct design implications: a forecasting model that is
mathematically superior but presented through a complex interface will
be rejected by SME owners in favour of a simpler, less accurate
alternative \cite{yuniarty2025}.

Somjaipeng and Jitsakul (2026) reached a complementary conclusion from
the auditing sector, applying the Technology-Organization-Environment
(TOE) framework and TAM to evaluate AI adoption among 120 SME-affiliated
auditors in Thailand. PLS-SEM analysis showed that environmental
factors --- particularly regulatory frameworks and industry norms ---
had a significant positive influence ($\beta = 0.562$) on perceived AI
benefits, explaining over 70\% of the variance in AI acceptance. The
implication for LocalDemand is that contextual fit (i.e., the app
looking and feeling like a tool designed specifically for food businesses)
is as important as technical performance \cite{somjaipeng2026}.

Aree and Ueasangkomsate (2025) studied 92 MSMEs in the Thai food supply
chain, confirming through linear regression analysis that AI adoption
significantly improves supply chain performance and firm traceability.
However, the study identified financial capital and technical know-how as
the two most significant barriers preventing MSMEs from adopting standard
AI solutions. The authors concluded that enterprise-grade AI is
effectively inaccessible to micro-businesses without substantial
subsidisation or simplification \cite{aree2025}.

Srivastava et al.\ (2026) reviewed how machine learning can be leveraged
by small and medium e-commerce enterprises for demand forecasting and
marketing analytics. Their primary contribution was identifying the
barriers: budget constraints, perceived model complexity, and inadequate
data infrastructure. The study argued that for ML to be viable for SMEs,
it must be delivered through accessible interfaces requiring zero
technical expertise \cite{srivastava2026}.

% ----------------------------------------------------------
\subsection{Mobile and Low-Cost Solutions}
% ----------------------------------------------------------

Two reviewed works directly address the mobile application delivery
layer. Sakuda, Leon, and Rojas (2024) developed an Android application
integrating the GPT-4 Large Language Model as an intelligent management
assistant for small shops in Peru. The app allowed owners to query
real-time sales data and receive conversational recommendations on
inventory and promotions. A qualitative survey revealed strong user
approval, with 66.67\% of surveyed shop owners rating the tool highly for
usability. The study demonstrates that mobile-first AI tools can be
successfully adopted by non-technical SME owners; however, the system's
reliance on generative AI (GPT-4) means its inventory recommendations
are based on contextual language generation rather than rigorous
time-series mathematics \cite{sakuda2024}.

Janchidfah et al.\ (2026) developed ``Stock Easy,'' a lightweight
inventory management mobile application evaluated using TAM with 132 retail
stakeholders. Their findings confirmed that perceived ease of use
(mean score = 4.52 out of 5) and process clarity were the strongest
predictors of user satisfaction, which in turn strongly correlated with
organisational readiness for digital transformation (Pearson r = 0.76).
The application, however, was purely reactive: it tracked current stock
levels manually and contained no predictive analytics. The study
explicitly identifies predictive demand forecasting as the next step
in the evolution of such tools --- the precise gap that LocalDemand
fills \cite{janchidfah2026}.

% ----------------------------------------------------------
\subsection{Generative AI vs.\ Specialised Forecasting Models}
% ----------------------------------------------------------

The emergence of large language models (LLMs) as potential supply chain
decision-support tools has prompted comparative evaluation. Vyakarnam
et al.\ (2025) systematically compared GPT-4.5, Grok 3, Gemini 2.0, and
DeepSeek R-1 against traditional ML methods (XGBoost, Random Forest) on
supply chain optimisation tasks for small businesses. While models such
as DeepSeek R-1 showed promise in synthesising demand planning insights
from retail datasets, all four generative models failed to correctly
solve deterministic combinatorial problems such as the Vehicle Routing
Problem. The authors concluded that LLMs are highly accessible and
cost-effective for narrative planning tasks but are fundamentally
ill-suited for rigorous numerical optimisation in supply chains
\cite{vyakarnam2025}.

This finding provides critical justification for LocalDemand's
architecture. Rather than using a generative LLM --- as Sakuda et al.\
(2024) did --- LocalDemand employs Facebook Prophet, a purpose-built
statistical forecasting engine with a well-defined mathematical
foundation. The trade-off is slightly higher technical complexity in the
back end, but the restock quantities generated are deterministic and
mathematically grounded, not probabilistic language outputs. The mobile
Flutter interface then absorbs all technical complexity on the front end,
delivering the ``ease of use'' that Yuniarty et al.\ (2025) and
Janchidfah et al.\ (2026) identify as the decisive adoption driver.

% ==========================================================
\section{Synthesis, Gaps, and Implications for LocalDemand}
% ==========================================================

% ----------------------------------------------------------
\subsection{Overall Strengths and Common Findings}
% ----------------------------------------------------------

The 25 works reviewed in this paper collectively establish several
consistent findings. First, machine learning and hybrid approaches
reliably outperform classical statistical models (ARIMA, SARIMA,
Holt-Winters) on food-sector demand forecasting tasks, with the
performance gap widening as dataset size and feature richness increase
\cite{gartner2024, ganguly2024, jatte2025}. Second, external regressors
--- particularly weather conditions, public holidays, and promotional
activities --- provide statistically significant improvements in forecast
accuracy across diverse contexts including restaurants, cafes, retail
chains, and cloud kitchens \cite{guler, kamboj2025, groene_zakharov,
lalitha2025, jain2024, qureshi2024}. Third, Facebook Prophet consistently
emerges as a pragmatic and well-validated middle ground for SME-facing
systems: it is competitive with LSTM in accuracy (85.71\% vs 92.31\%
in Jatte et al.; RMSE 0.46 vs 0.1535 in Aqila et al.), substantially
easier to configure and maintain, and natively handles missing data,
trend shifts, holiday effects, and custom regressors \cite{jatte2025,
aqila2026, kamboj2025, guler}. Fourth, technology adoption research
consistently finds that perceived ease of use and accessibility are
more decisive than raw model performance for SME operators, making the
delivery interface as important as the underlying algorithm
\cite{yuniarty2025, janchidfah2026, aree2025}.

% ----------------------------------------------------------
\subsection{Major Research Gaps}
% ----------------------------------------------------------

Despite this rich body of work, four interconnected gaps remain
unaddressed in the existing literature:

\textbf{Gap 1 --- Forecast-to-Action Gap:} The vast majority of reviewed
forecasting studies --- including Aqila et al.\ (2026), Jatte et al.\
(2025), Sandeep et al.\ (2025), Ganguly \& Mukherjee (2024), and Gartner
et al.\ (2024) --- terminate their analysis at error metrics (RMSE, MAPE,
MAE). They do not convert predictions into actionable inventory decisions
such as restock quantities or purchase orders. The few studies that do
attempt this step (Rahmatillah et al.\ 2026; Munawar \& Sudirman 2026)
rely on data volumes (nearly 1{,}000 days of history) that are
unavailable to newly established food businesses \cite{aqila2026,
jatte2025, rahmatillah2026, munawar2026}.

\textbf{Gap 2 --- Accessibility and Mobile Integration Gap:} No reviewed
work combines a rigorous time-series forecasting engine with a
mobile-first delivery layer specifically designed for non-technical SME
operators. Sakuda et al.\ (2024) provide a mobile app but use GPT-4
rather than a deterministic forecasting model \cite{sakuda2024}.
Janchidfah et al.\ (2026) provide a mobile inventory app but with no
predictive capability \cite{janchidfah2026}. The gap between ``accurate
but complex'' and ``simple but not forecasting'' is unoccupied.

\textbf{Gap 3 --- Micro-SME and Food-Sector Specificity Gap:} Several
high-performing systems (AID-PAF, Smart Stock, GRU supply chain model)
require hardware infrastructure, large datasets, or computational
resources that are incompatible with single-location micro-SMEs such as
a neighbourhood bakery or a small independent café \cite{aidpaf,
smartstock, qureshi2024}. Research that combines high model accuracy with
SME-scale constraints simultaneously is scarce.

\textbf{Gap 4 --- Cold-Start and Sparse-Data Gap:} As established in
Section~IV, most deep learning models evaluated in the literature require
several hundred to thousands of training samples to converge reliably.
This makes them unsuitable for new businesses or new product lines, a
scenario that is ubiquitous in the food SME sector. No reviewed work
provides a complete solution that addresses sparse data, external
regressor integration, and operational output generation
simultaneously \cite{susanto2026, gartner2024, rahmatillah2026}.

% ----------------------------------------------------------
\subsection{How LocalDemand Fills These Gaps}
% ----------------------------------------------------------

LocalDemand is designed as a direct response to all four identified gaps,
with each architectural decision traceable to the evidence reviewed above.

Against \textbf{Gap 1}, LocalDemand does not terminate at forecast
generation. Following the simulation methodology of Rahmatillah et al.\
(2026), it converts Prophet's daily demand forecasts into concrete restock
recommendations --- expressed as units-to-order per ingredient or menu
item --- with a configurable safety-stock buffer. This makes the system's
primary output an operational decision, not a statistical metric.

Against \textbf{Gap 2}, LocalDemand delivers its forecasts and restock
recommendations through a Flutter mobile application, making it accessible
on any Android or iOS device without installation of data-science
tooling. In contrast to Sakuda et al.'s (2024) LLM-based app, the
recommendations in LocalDemand are generated by Facebook Prophet --- a
deterministic, mathematically grounded time-series model --- ensuring
that inventory quantities are precise and reproducible. The interface
design is governed by the TAM findings of Yuniarty et al.\ (2025) and
Janchidfah et al.\ (2026): simplicity, clarity, and friction-free access
are non-negotiable design requirements.

Against \textbf{Gap 3}, LocalDemand is explicitly scoped to
single-location food micro-SMEs (cafes, bakeries, restaurants) operating
with minimal data infrastructure. It requires only a smartphone and basic
daily sales records --- no IoT devices, no dedicated servers, and no
data-science team. This positions it in the micro-SME accessibility space
validated by Aree \& Ueasangkomsate (2025) and Srivastava et al.\ (2026)
as urgently needed but currently unserved.

Against \textbf{Gap 4}, LocalDemand selects Facebook Prophet specifically
for its cold-start resilience. Prophet's Bayesian prior structure and
Fourier-series seasonality estimation allow it to produce reasonable
forecasts from as few as a few weeks of daily observations, and it handles
zero-valued days and missing entries natively --- directly addressing the
zero-inflation problem documented by Susanto et al.\ (2026). The addition
of weather and holiday regressors provides independent predictive signal
that stabilises forecasts even when internal sales history is sparse, as
demonstrated by Guler et al., Kamboj et al., and Groene \& Zakharov
\cite{guler, kamboj2025, groene_zakharov, susanto2026}.

In summary, LocalDemand occupies a design space that the existing
literature has identified as both urgently needed and currently vacant:
a system that is simultaneously accurate (Prophet + weather/holiday
regressors), actionable (daily restock quantities with safety buffer),
accessible (Flutter mobile app), and robust to micro-SME data realities
(cold-start and sparse-data resilience).

\begin{table*}[t]
\centering
\caption{Comparative Analysis of Demand Forecasting Approaches in Food and Retail SMEs}
\label{tab:comparative_analysis}
\resizebox{\textwidth}{!}{%
\renewcommand{\arraystretch}{1.4}
\begin{tabular}{@{}p{3.5cm}p{2.5cm}p{3cm}p{3cm}p{3.5cm}p{2cm}p{1.5cm}p{4cm}@{}}
\hline
\textbf{Study / Authors (Year)} & \textbf{Model(s) Used} & \textbf{Key External Regressors} & \textbf{Dataset Size / Context} & \textbf{Performance Metrics} & \textbf{Actionability} & \textbf{SME Suitability} & \textbf{Main Limitation / Gap Addressed} \\ \hline

Groene \& Zakharov & XGBoost & Weather, Holidays, Events & 5+ years hourly F\&B sales & 22--33\% RMSE reduction & Forecast-only & Medium & Lacks operational restock outputs \\

\textbf{Guler et al.} & \textbf{Facebook Prophet} & \textbf{Weather, Special Days} & \textbf{Restaurant daily data} & \textbf{MAE = 0.76, RMSE = 0.82} & \textbf{Forecast-only} & \textbf{High} & \textbf{Lacks restock simulation / mobile app} \\

Jatte et al. (2025) & LSTM vs \textbf{Prophet} & None specified & Large retail dataset & LSTM 92\% acc, Prophet 85.7\% & Forecast-only & Medium & \textbf{Prophet} highlighted for low overhead \\

Aqila et al. (2026) & LSTM, \textbf{Prophet}, SARIMA & None specified & Hourly coffee shop sales & LSTM RMSE 0.15, Prophet 0.46 & Forecast-only & Medium & Stops at metrics; no restock logic \\

\textbf{Kamboj et al. (2025)} & \textbf{Prophet + XGBoost} & \textbf{Promotions, Holidays} & \textbf{State-level retail} & \textbf{12--18\% MAPE reduction} & \textbf{Forecast-only} & \textbf{Medium} & \textbf{Prophet} handles holidays robustly \\

Rahmatillah et al. (2026) & LSTM vs Moving Avg & None specified & SME ice cream ($\sim$1k days) & RMSE 1.65, zero stockouts & Simulation & Medium & High data requirement (cold-start gap) \\

Munawar \& Sudirman (2026) & LSTM & None specified & Hourly SME F\&B data & Lowest MAE/RMSE & Simulation & Low & Excludes weather/holidays; high data need \\

Gartner et al. (2024) & CNN, ARIMA, GAM & None specified & SME real-world data & Normalised RMSE 0.31 & Forecast-only & Medium & Deep learning fails on cold-start \\

Lalitha et al. (2025) & Random Forest, ARIMA & Weather, Cost data & Perishable goods sales & MSE 53.97 & Pricing Engine & Low & Needs continuous real-time data feeds \\

AID-PAF Framework & TFNA (Deep Learning) & Weather, Holidays, Sentiment & Simulated digital twin & MAPE 6.5\%, 40\% waste red. & Optimization & Low & Hardware/computing demands too high \\

Smart Stock & XGBoost + IoT & None specified & Small warehouse & 96\% classification accuracy & Alerts & Low & Hardware dependency (IoT/Edge) \\

Susanto et al. (2026) & XGBoost vs MA(7) & None specified & Zero-inflated SME data & MA(7) outperformed ML & Forecast-only & Medium & Regression fails on intermittent data \\ \hline

\textbf{LocalDemand (Proposed)} & \textbf{Facebook Prophet} & \textbf{Weather, Holidays} & \textbf{Micro-SME daily sales} & \textbf{TBD (Evaluation phase)} & \textbf{Restock Recs.} & \textbf{High} & \textbf{Fills forecast-to-action \& mobile gaps} \\ \hline
\end{tabular}%
}
\end{table*}


% ==========================================================
\section{Conclusion}
% ==========================================================

This literature review has synthesized 25 research works spanning
demand forecasting methods, external factor integration, SME-specific
data challenges, inventory optimisation, and technology adoption. The
evidence collectively supports the design rationale of LocalDemand across
five dimensions.

In terms of \textbf{model selection}, Facebook Prophet is empirically
validated as a competitive, interpretable, and operationally practical
forecasting engine for food-sector SMEs, outperforming classical
statistical models and approaching deep-learning accuracy while requiring
a fraction of the data and computational infrastructure
\cite{jatte2025, aqila2026, kamboj2025, guler}.

In terms of \textbf{external regressors}, weather and holiday data are
the two most consistently evidence-backed demand drivers across the
reviewed food-sector literature. Their native support in Prophet's
regressor framework makes their integration straightforward and well
precedented \cite{guler, kamboj2025, groene_zakharov, lalitha2025,
jain2024}.

In terms of \textbf{actionability}, the simulation studies of Rahmatillah
et al.\ (2026) and Munawar \& Sudirman (2026) establish that translating
forecasts into purchasing decisions eliminates stockouts and reduces
total inventory costs, validating the restock-recommendation layer of
LocalDemand \cite{rahmatillah2026, munawar2026}.

In terms of \textbf{accessibility}, the TAM and PLS-SEM studies of
Yuniarty et al.\ (2025), Somjaipeng \& Jitsakul (2026), Aree \&
Ueasangkomsate (2025), and the mobile-app usability findings of
Janchidfah et al.\ (2026) and Sakuda et al.\ (2024) confirm that
a mobile-first, low-friction delivery interface is a necessary condition
--- not merely a desirable feature --- for adoption by micro-SME
operators \cite{yuniarty2025, somjaipeng2026, aree2025,
janchidfah2026, sakuda2024}.

In terms of \textbf{data realism}, the cold-start and sparse-data
findings of Susanto et al.\ (2026), Gartner et al.\ (2024), and
Rahmatillah et al.\ (2026) confirm that the dominant deep learning
methods in the forecasting literature are incompatible with the data
realities of newly established or small-scale food businesses, and
that Prophet-based approaches with external regressors represent the
most viable and well-supported alternative \cite{susanto2026,
gartner2024, rahmatillah2026}.

Taken together, the 25 reviewed works provide a coherent and complete
evidence base for the architecture of LocalDemand: Facebook Prophet
as the forecasting core, weather and holiday data as primary external
regressors, a restock-recommendation layer with a safety-stock buffer
as the operational output, and a Flutter mobile application as the
accessible, user-facing delivery channel. Future research directions
identified by the reviewed literature include the integration of social
media sentiment signals \cite{aidpaf}, multi-outlet and cross-location
demand modelling, and longitudinal evaluation of inventory cost
reduction in real micro-SME deployments.

% ==========================================================
%  BIBLIOGRAPHY
%  (Add your actual .bib file or replace with thebibliography
%   entries as needed for your institution's submission format)
% ==========================================================

\begin{thebibliography}{99}

\bibitem{aqila2026}
K. S. Aqila et al., ``Comparative Evaluation of SARIMA, FB Prophet, and
LSTM for Hourly Demand Forecasting in Saudagar Kopi,'' 2026.

\bibitem{sandeep2025}
Sandeep V et al., ``Smart Sales Forecasting Machine Learning Models for
Demand Prediction in Retail,'' 2025.

\bibitem{kamboj2025}
M. Kamboj et al., ``State-Level Drug Retail Sales Forecasting Using
XGBoost and Facebook Prophet,'' 2025.

\bibitem{gartner2024}
T. Gartner et al., ``Automated Demand Forecasting in Small to
Medium-Sized Enterprises,'' 2024.

\bibitem{ganguly2024}
P. Ganguly and I. Mukherjee, ``Enhancing Retail Sales Forecasting with
Optimized Machine Learning Models,'' 2024.

\bibitem{jatte2025}
H. Jatte et al., ``Optimization of Forecasting Performance in the Retail
Sector Using Artificial Intelligence,'' 2025.

\bibitem{groene_zakharov}
Groene and Zakharov, ``AI-Based Sales Forecasting for F\&B Businesses.''

\bibitem{guler}
Guler et al., ``Forecasting Restaurant Sales with Weather and Special
Days using Facebook Prophet.''

\bibitem{susanto2026}
J. C. Susanto et al., ``Performance Analysis of Time Series Models to
Predict Sales Activity of SMEs,'' 2026.

\bibitem{munawar2026}
M. A. Munawar and I. D. Sudirman, ``LSTM-Based Hourly Sales Forecasting
in SME Food and Beverage Retail,'' 2026.

\bibitem{vyakarnam2025}
S. Vyakarnam et al., ``Optimizing Supply Chain Management for Small and
Medium-Sized Businesses Using AI,'' 2025.

\bibitem{yuniarty2025}
Yuniarty et al., ``Technology Acceptance and Security Influences on AI
Adoption among Small and Medium Enterprises,'' 2025.

\bibitem{aree2025}
T. Aree and P. Ueasangkomsate, ``Technology Adoption of MSMEs in Food
Supply Chain: The Role of AI in Traceability and Transparency,'' 2025.

\bibitem{somjaipeng2026}
K. Somjaipeng and W. Jitsakul, ``Factors influencing audit quality
through the acceptance and use of AI by auditors in Thailand,'' 2026.

\bibitem{srivastava2026}
S. Srivastava et al., ``Predictive Marketing and Sales Analytics Using
Machine Learning for Small and Medium E-Commerce Enterprises,'' 2026.

\bibitem{lalitha2025}
Lalitha R et al., ``Enhanced Sales Forecasting of Perishable and
Non-Perishable Goods Using Weather Data with Insights on Imports and
Price Sensitivity,'' 2025.

\bibitem{jain2024}
N. Jain et al., ``Weather-Integrated Recommendations for Food Choices
with SHAP Explanation,'' 2024.

\bibitem{qureshi2024}
N. U. H. Qureshi et al., ``Demand Forecasting in Supply Chain Management
for Rossmann Stores using Weather Enhanced Deep Learning Model,'' 2024.

\bibitem{rahmatillah2026}
I. Rahmatillah, Herwanto, and Sudirman, ``Evaluation of Optimized LSTM
and Moving Average in SME Inventory Forecasting and Simulation,'' 2026.

\bibitem{smallbiz_deepar}
Author(s) unspecified, ``Demand Forecasting System, Small Business
Optimization.''

\bibitem{aidpaf}
Author(s) unspecified, ``AI Predictive Analytics for Sustainable
Restaurant Operations (AID-PAF).''

\bibitem{smartstock}
Author(s) unspecified, ``Smart Stock: Small Warehouse Inventory.''

\bibitem{foodwaste_ml}
Author(s) unspecified, ``Predictive Analytics, Food Supply Chain and
Waste.''

\bibitem{sakuda2024}
G. Sakuda, Leon, and Rojas, ``Mobile Application for the Intelligent
Management of Small Shops Using the GPT-4 Model,'' 2024.

\bibitem{janchidfah2026}
T. Janchidfah et al., ``A Mobile Application for Inventory Management
to Support Digital Transformation in Small Retail Businesses,'' 2026.

\end{thebibliography}

\end{document}
